\documentclass[11pt]{amsart}

\usepackage[T1]{fontenc}
\usepackage{lmodern}
\usepackage{amsmath,amssymb,amsthm,mathtools,mathrsfs}
\usepackage{enumitem}
\usepackage[a4paper,margin=32mm]{geometry}
\usepackage[hidelinks]{hyperref}

\newcommand{\C}{\mathbb C}
\newcommand{\Z}{\mathbb Z}
\newcommand{\Q}{\mathbb Q}
\renewcommand{\P}{\mathbb P}
\newcommand{\cO}{\mathcal O}
\newcommand{\cX}{\mathcal X}
\newcommand{\cS}{\mathcal S}
\newcommand{\cT}{\mathcal T}
\newcommand{\cK}{\mathcal K}
\newcommand{\cD}{\mathcal D}
\newcommand{\Sing}{\operatorname{Sing}}
\newcommand{\Pic}{\operatorname{Pic}}
\newcommand{\Spec}{\operatorname{Spec}}
\newcommand{\id}{\operatorname{id}}

\theoremstyle{plain}
\newtheorem{theorem}{Theorem}[section]
\newtheorem{proposition}[theorem]{Proposition}
\newtheorem{lemma}[theorem]{Lemma}
\newtheorem{corollary}[theorem]{Corollary}

\theoremstyle{definition}
\newtheorem{definition}[theorem]{Definition}
\newtheorem{example}[theorem]{Example}

\theoremstyle{remark}
\newtheorem{remark}[theorem]{Remark}

\title[An independent test family]{The Singular-Fibre Vanishing Step: An Independent Audit and Proper Test Family}
\author{Human author to be supplied}
\date{1 September 2026}

\begin{document}

\begin{abstract}
We isolate an endpoint implication that is not justified by the displayed
sheaf sequence in the singular-fibre argument of CDP20 (p.~692).
Levent Alp{\"o}ge's supplied source
manuscript already states and proves, for its proposed fibre, a
conductor-supported differential and uses it to challenge that reduction
\cite[Sec.~10.2, Lemma~10.3, source PDF pp.~87--88]{S6Source}; we neither
claim that candidate-specific observation nor repeat its $S^6$ construction.
The independent calculation here is a proper
test family.
For a reduced effective
Cartier divisor $S$ in a smooth complex threefold $X$, the conormal sequence
has to be filtered through the torsion subsheaf of $\Omega^1_S$.  We prove
that, after the two endpoint vanishings used in CDP20, the ambient section
space is the kernel of an explicit connecting homomorphism
\[
\delta_A:H^0(S,\mathscr T_S\otimes A)\longrightarrow
H^1(S,N^*_{S/X}\otimes A),
\]
not automatically zero.  We then give a proper one-parameter morphism over a
disc with smooth total space whose central fibre is a nodal cubic times
$\P^1$.  A
non-torsion gluing line bundle on that fibre has both endpoint groups zero,
while the ambient twisted cotangent space is nonzero.  Thus the formal
  three-sheaf implication used at p.~692 of CDP20, if read for an arbitrary
  fibre line bundle, is not a formal consequence of the displayed sequence
  without a hypothesis controlling the resulting connecting map; extendability of a
  line bundle alone is not such a control.  The example can in fact be chosen
  with an extending line bundle on the displayed local total space, but it
  does not satisfy the compact global hypotheses of CDP20 Theorem~2.2 and does
  not claim a
counterexample to that theorem.  We also give a derived reformulation of the
surviving CDP20 criterion: one twist with $H^2(X,T_X\otimes L_0)=0$ is enough
to recover the
generic vanishing and $c_3\leq0$ conclusions, without using the defective
endpoint reduction.  Finally, on the bare local models used by the proposed
construction, we compute the ordinary relative one-form module: it is
torsion-free but non-locally-free at a crossing, whereas a multiplicity-$m$
  smooth cyclic chart has torsion exactly $R/(s^{m-1})\,ds$.  This is a local correction
for $p=1$, not a global assertion about the proposed threefold.  The
  split-extension coinvariant calculation already appears in the source
  manuscript (Sec.~10.5, Lemma~10.7, source PDF pp.~90--91) and is deliberately
  not repeated here.
\end{abstract}

\maketitle

\section{Purpose and precise scope}

Campana, Demailly, and Peternell prove a vanishing theorem for compact
complex threefolds with vanishing first and second integral cohomology and
algebraic dimension one \cite{CDP2020}.  In the singular-fibre part of their
proof they write, for an irreducible reduced fibre $S$, the conormal sequence
and reduce the desired vanishing to
\[
 H^0(S,\widetilde\Omega^1_S\otimes L|_S)=0,
 \qquad
 \widetilde\Omega^1_S:=\Omega^1_S/\text{torsion}.
 \tag{1.1}
\]
The purpose of this note is narrowly defined: we make the omitted connecting
map explicit, give a proper geometric example showing that the two endpoint
vanishings are not by themselves sufficient, and record a derived global
criterion that isolates the extra vanishing needed by the intact part of the
CDP strategy.  The earlier version of this strategy is the 1998 paper
\cite{CDP1998}; the present note addresses the singular-fibre reduction in the
  later paper.  Levent Alp{\"o}ge's source manuscript, Section~10, gives the corresponding
  candidate-specific normalization--conductor argument (Sec.~10.2,
  Lemma~10.3, source PDF pp.~87--88; the public audit deposit is
  \cite{S6AuditArchive}); the filtration and conductor descent are recalled
  here only to state an independent test precisely.  No priority is claimed
  for them.  We do not infer
that CDP20 Theorem~2.2 is
false, nor do we use the example to assert the existence of a complex
structure on $S^6$.

All sheaves below are holomorphic sheaves.  If $S$ is reduced, let
$j:S_{\rm reg}\hookrightarrow S$ be the inclusion of the smooth locus and
put
\[
 \mathscr T_S:=\ker\bigl(\Omega^1_S\longrightarrow
 j_*\Omega^1_{S_{\rm reg}}\bigr),
 \qquad
 \widetilde\Omega^1_S:=\Omega^1_S/\mathscr T_S.
 \tag{1.2}
\]
For the reduced Cartier divisors used below, $\mathscr T_S$ is precisely the
torsion subsheaf of $\Omega^1_S$.

\begin{lemma}[identification of the torsion subsheaf]
\label{lem:torsion-identification}
If $S$ is a reduced effective Cartier divisor in a smooth complex manifold,
then
\[
\ker\bigl(\Omega^1_S\longrightarrow j_*\Omega^1_{S_{\rm reg}}\bigr)
  =\operatorname{Tor}(\Omega^1_S).
\]
\end{lemma}

\begin{proof}
Every irreducible component of a reduced hypersurface is generically smooth,
so $S_{\rm reg}$ is dense and its complement is a proper analytic subset.
If a local section of $\Omega^1_S$ vanishes on $S_{\rm reg}$, its support avoids
all minimal primes of the reduced local ring.  Prime avoidance therefore gives
an element of its annihilator outside every minimal prime; that element is a
nonzerodivisor, so the section is torsion.  Conversely,
$\Omega^1_S|_{S_{\rm reg}}$ is locally free, hence has no torsion, and every
torsion section restricts to zero there.
\end{proof}

\section{The exact cotangent filtration}

Let $i:S\hookrightarrow X$ be a reduced effective Cartier divisor in a
smooth complex threefold.  Write
\[
 N^*_{S/X}:=\mathcal I_S/\mathcal I_S^2,
 \qquad
 \Omega^1_X|_S:=i^*\Omega^1_X.
\]

\begin{lemma}[left exactness of the conormal sequence]
\label{lem:conormal}
There is a short exact sequence
\[
0\longrightarrow N^*_{S/X}\longrightarrow \Omega^1_X|_S
\xrightarrow{\rho_0}\Omega^1_S\longrightarrow0.
\tag{2.1}
\]
\end{lemma}

\begin{proof}
Right exactness is the conormal sequence for a closed immersion.  We check
the left end.  Work at a point of $S$ and let $R$ be the regular local ring
of $X$, with reduced local equation $g$ for $S$.  The map at the left is
\[
 R/(g)\longrightarrow (\Omega^1_R/g\Omega^1_R),
 \qquad \bar h\longmapsto h\,dg.
 \tag{2.2}
\]
Let $\mathfrak p$ be a minimal prime of $(g)$.  Since $R$ is regular and
has characteristic zero, the height-one local ring $R_{\mathfrak p}$ is a
discrete valuation ring.  Reducedness of $(g)$ gives
$\operatorname{ord}_{\mathfrak p}(g)=1$, so $g$ is a unit times a
uniformizer there; in characteristic zero the differential of that
uniformizer is nonzero modulo $\mathfrak p$.  Hence $dg$ is nonzero in
$\Omega^1_{R_{\mathfrak p}}\otimes k(\mathfrak p)$, so (2.2) sends $\bar h$
to zero only if $h\in\mathfrak p$.  This holds for every minimal prime.
Because $R/(g)$ is reduced, the intersection of its minimal primes is zero;
therefore $\bar h=0$.  Thus (2.1) is exact.
\end{proof}

Define the inverse-image sheaf
\[
 \mathscr K_S:=\ker\bigl(\Omega^1_X|_S\xrightarrow{\rho_0}
 \Omega^1_S\longrightarrow\widetilde\Omega^1_S\bigr).
 \tag{2.3}
\]

\begin{proposition}[the missing filtration]
\label{prop:filtration}
There are natural short exact sequences
\[
0\longrightarrow N^*_{S/X}\longrightarrow\mathscr K_S
\longrightarrow\mathscr T_S\longrightarrow0,
\tag{2.4}
\]
\[
0\longrightarrow\mathscr K_S\longrightarrow\Omega^1_X|_S
\longrightarrow\widetilde\Omega^1_S\longrightarrow0.
\tag{2.5}
\]
\end{proposition}

\begin{proof}
By definition, $\mathscr K_S$ is the inverse image of $\mathscr T_S$ under
the surjection $\rho_0$ in (2.1).  Restricting $\rho_0$ to this inverse
image gives a surjection $\mathscr K_S\twoheadrightarrow\mathscr T_S$;
its kernel is $N^*_{S/X}$ by (2.1).  This proves (2.4).  The quotient map
in (2.5) is surjective because $\rho_0$ is surjective and
$\Omega^1_S\twoheadrightarrow\widetilde\Omega^1_S$ is the quotient by
$\mathscr T_S$.
\end{proof}

\begin{proposition}[exact global-section criterion]
\label{prop:criterion}
Let $A$ be a line bundle on $S$.  If
\[
 H^0(S,\widetilde\Omega^1_S\otimes A)=0,
 \tag{2.6}
\]
then the inclusion in (2.5) induces an isomorphism
\[
 H^0(S,\Omega^1_X|_S\otimes A)
 \simeq H^0(S,\mathscr K_S\otimes A).
 \tag{2.7}
\]
If in addition
\[
 H^0(S,N^*_{S/X}\otimes A)=0,
 \tag{2.8}
\]
then there is a canonical isomorphism
\[
 H^0(S,\Omega^1_X|_S\otimes A)
 \simeq \ker\!\left(
 H^0(S,\mathscr T_S\otimes A)
 \xrightarrow{\;\delta_A\;}
 H^1(S,N^*_{S/X}\otimes A)
 \right),
 \tag{2.9}
\]
where $\delta_A$ is the connecting map of (2.4) tensored by $A$.
Consequently, under (2.6) and (2.8), the ambient group vanishes if and only
if $\delta_A$ is injective.
\end{proposition}

\begin{proof}
Tensor (2.5) by the locally free sheaf $A$ and take global sections.  If
(2.6) holds, the long exact sequence gives (2.7).  Tensor (2.4) by $A$ and
take its first four terms.  Condition (2.8) identifies the middle global
sections with the kernel of the connecting map.  Combining the two
identifications gives (2.9).
\end{proof}

\begin{remark}
The map $\delta_A$ is not an auxiliary hypothesis introduced for
terminological convenience: it is the exact obstruction to lifting a torsion
section of $\Omega^1_S\otimes A$ through the conormal extension.  In
\v{C}ech terms, choose local lifts of a section to
$\mathscr K_S\otimes A$; their pairwise differences are sections of
$N^*_{S/X}\otimes A$ and represent $\delta_A$.  The two endpoint groups in
(2.6) and (2.8) contain no information about this obstruction.
\end{remark}

\section{A conductor-supported section}

The normalization--conductor mechanism in this section is a self-contained
restatement (and slight sheaf-theoretic generalization) of the source
manuscript's Lemma~10.3; it is included so that the independent test family
can be checked without importing an unprinted step, and no priority is
claimed for it.  We use it to produce a nonzero element of the kernel in
(2.9).  Write
\(\theta_g:=i^*(dg)\in H^0(S,\Omega^1_X|_S)\).  Whenever we write
$dg|_S$ below it means this ambient restriction $i^*(dg)$, not the intrinsic
differential $d(g|_S)=0$ of the constant function on the fibre.  Suppose
$g:X\to\Delta$ is a holomorphic map from a smooth complex threefold to a
disc, $S=g^{-1}(0)$ is reduced, and, writing $D=\Sing(S)$, every point of
$D$ has a normal-crossing (self-crossing) map-germ model: in local coordinates
$(u,v,w)$, $g=h\,uv$ with $h$ a holomorphic unit.  After replacing $u$ by
$hu$, we use the equivalent model
\[
 (X,g,S)\simeq\bigl(\C^3,uv,\{uv=0\}\bigr).
 \tag{3.1}
\]
At points of $S\setminus D$ use the smooth model $(\C^3,\{u=0\})$.
Let $\eta:\widetilde S\to S$ be the normalization.
For the application below, $D$ is smooth and $\eta^{-1}(D)$ is the disjoint
union of two copies of $D$.

\begin{lemma}[normalization--conductor sequence]
\label{lem:conductor}
In the situation above there is an exact sequence
\[
0\longrightarrow\cO_S\longrightarrow
 \eta_*\cO_{\widetilde S}\oplus\cO_D
\longrightarrow\eta_*\cO_{\eta^{-1}(D)}\longrightarrow0,
\tag{3.2}
\]
where the maps are
\[
h\longmapsto(\eta^*h,h|_D),\qquad
(\widetilde h,k)\longmapsto\widetilde h|_{\eta^{-1}(D)}-\eta^*k.
\]
\end{lemma}

\begin{proof}
The assertion is local.  At a point of $D$ the local ring of $S$ is
\(R=\C\{u,v,w\}/(uv)\).  Its normalization is
\(\C\{u,w\}\oplus\C\{v,w\}\), the conductor is the $w$-axis, and the
last term in (3.2) is two copies of $\C\{w\}$.  The resulting map is
\[
 (a(u,w),b(v,w),c(w))\longmapsto
 (a(0,w)-c(w),b(0,w)-c(w)).
 \tag{3.3}
\]
Its kernel consists exactly of pairs obtained by restricting a function on
the union $uv=0$; for example, a kernel element is represented by
$a(u,w)+b(v,w)-c(w)$.  Surjectivity follows by taking $a$ and $b$ to be
the two prescribed functions and $c=0$.  Away from $D$ the normalization is
an isomorphism, so these local calculations prove exactness.
\end{proof}

\begin{proposition}[conductor differential]
\label{prop:conductor-differential}
Let $A$ be a line bundle on $S$ such that $\eta^*A\simeq\cO_{\widetilde S}$,
and choose a nowhere-vanishing trivializing section $e$ of $\eta^*A$.
Put $E=\Omega^1_X|_S\otimes A$.  Assume that $\eta^*\theta_g$ restricts to
zero on $\eta^{-1}(D)$.  Then there is a unique section
\[
 0\ne s\in H^0(S,\Omega^1_X|_S\otimes A)
 \tag{3.4}
\]
whose image in the normalization term of (3.2), namely
$\eta_*(\eta^*E)\oplus E|_D$, is
\[
  \bigl(\eta^*\theta_g\otimes e,\,0\bigr).
 \tag{3.5}
\]
The image of $s$ in $\Omega^1_S\otimes A$ is supported on $D$, hence is a
torsion section and maps to zero in $\widetilde\Omega^1_S\otimes A$.
\end{proposition}

\begin{proof}
Tensor (3.2) by the locally free sheaf
$E=\Omega^1_X|_S\otimes A$.  Finiteness of $\eta$ and the projection formula
identify the first and third terms with $E$ and
$\eta_*(\eta^*E|_{\eta^{-1}(D)})$.  The two components in (3.5) have the
same restriction to $\eta^{-1}(D)$ because the first one is zero there.
Exactness therefore gives a unique section $s$.

On $S\setminus D$, the normalization is an isomorphism and $s$ is
$dg|_S\otimes e$.  Since $g$ is a submersion along the smooth fibre, this is
nonzero on a dense open set, so $s\ne0$.  Its intrinsic image vanishes on
$S\setminus D$ because $g$ is constant on $S$ there.  To see explicitly that
this makes the image torsion, work at a point $x\in S$.  The support of the
germ is contained in the proper closed set $D$; hence, for the finitely
generated module generated by the germ, a power of the ideal of $D$ kills it.
Because $S$ is reduced and no irreducible component of $S$ is contained in
$D$, prime avoidance supplies an element of that ideal outside every minimal
prime of $\mathcal O_{S,x}$.  A suitable power is therefore a non-zero-divisor
annihilating the germ.  The image is torsion and is killed in the quotient by
$\mathscr T_S$.
\end{proof}

\begin{corollary}
\label{cor:normal-crossing-kernel}
If, in addition, $H^0(S,N^*_{S/X}\otimes A)=0$ and
$H^0(S,\widetilde\Omega^1_S\otimes A)=0$, then the section in
Proposition~\ref{prop:conductor-differential} proves that $\delta_A$ in
(2.9) is not injective.
\end{corollary}

\begin{proof}
The section is nonzero and maps to zero in the quotient
$\widetilde\Omega^1_S=\Omega^1_S/\mathscr T_S$, so it
lies in the left-hand side of (2.9).  Proposition~\ref{prop:criterion}
identifies that group with $\ker\delta_A$, which is therefore nonzero.
\end{proof}

\section{The proper nodal-cubic test family}

We now instantiate the preceding mechanism in an explicit proper family.
Let $\Delta$ be a sufficiently small disc with coordinate $t$, and define
\[
 \mathscr C=
 \left\{([x:y:z],t)\in\P^2\times\Delta:
 y^2z=x^2(x+z)+tz^3\right\}.
 \tag{4.1}
\]
Let $g:\mathcal X\to\Delta$ be the projection from
$\mathcal X=\mathscr C\times\P^1$, and put
\[
 C_0=\mathscr C\cap\{t=0\},\qquad S=C_0\times\P^1=g^{-1}(0).
 \tag{4.2}
\]

\begin{proposition}[geometry of the family]
\label{prop:family-geometry}
The projection $g$ is proper, $\mathcal X$ is a smooth complex threefold,
  and $S$ is a reduced irreducible normal-crossing (self-crossing) surface whose normalization
is $\P^1\times\P^1$.  Its singular locus is
\(D=\{[0:0:1]\}\times\P^1\), and locally along $D$ the pair $(\mathcal X,S)$
has the model (3.1); away from $D$ the fibre is smooth and has the usual
local model $(\C^3,\{u=0\})$.
\end{proposition}

\begin{proof}
Properness follows because $\mathscr C$ is closed in the projective
projection $\P^2\times\Delta\to\Delta$.  Write the defining equation as
\[
 F=y^2z-x^2(x+z)-tz^3.
\]
On the chart $z\ne0$, $\partial F/\partial t=-z^3\ne0$, so $\mathscr C$
is smooth there.  On $z=0$, the equation forces $x=0$ and $y\ne0$; at such
points $\partial F/\partial z=y^2\ne0$.  Hence $\mathscr C$, and therefore
$\mathcal X$, is smooth.  The only finite discriminant values of the fibres
are $t=0$ and $t=-4/27$; after shrinking $\Delta$ to $|t|<4/27$, every
noncentral fibre is smooth.

The central curve is
\(C_0:\ y^2z=x^2(x+z)\).  Its only singular point is
\(n=[0:0:1]\): on $z=1$ the equations $F=F_x=F_y=0$ give $(x,y)=(0,0)$,
and the point at infinity $[0:1:0]$ is smooth because $F_z=1$ there
(Euler's identity also gives the remaining projective singularity equation).
The parametrization
\[
 \nu:\P^1\longrightarrow C_0,\qquad
 [u:v]\longmapsto
 [(u^2-v^2)v:\,u(u^2-v^2):\,v^3]
 \tag{4.3}
\]
maps the two points $p=[1:1]$, $q=[-1:1]$ to $n$.  On the chart $v\ne0$, set
$r=u/v$; after scaling, its image is
$[r^2-1:r(r^2-1):1]$.  Conversely, for an affine point of $C_0$ with
$x\ne0$, the equation gives $r=y/x$ and $x=r^2-1$, so it is hit by
$[r:1]$.  The only affine point with $x=0$ is the node $(0,0)$, hit by
$[1:1]$ and $[-1:1]$, and the point at infinity $[0:1:0]$ is hit by
$[1:0]$.  Thus $\nu(\P^1)=C_0$; as a nonconstant morphism to a projective
curve it is finite and generically one-to-one, hence the normalization.  In
particular $C_0$ is irreducible and rational.  The product normalization
is $\eta=\nu\times\id_{\P^1}:\P^1\times\P^1\to S$.

At $n$ the affine equation is
\(t=y^2-x^2-x^3\).  On a sufficiently small neighbourhood choose the
holomorphic square root $h(x)=\sqrt{1+x}$ and set
\(u=y-xh(x)\), \(v=y+xh(x)\).  Then $uv=t$, and the map
$(x,y)\mapsto(u,v)$ has invertible $2\times2$ Jacobian at $(0,0)$, so $u,v$
are local coordinates on the surface $\mathscr C$.  If $s$ is a local
coordinate on $\P^1$, then $(u,v,s)$ are coordinates on $\mathcal X$ and
$t=uv$.  Hence the pair has the asserted normal-crossing (self-crossing)
local model, with singular locus $D$ as stated.
\end{proof}

\begin{lemma}[the non-torsion gluing bundle]
\label{lem:gluing-bundle}
Fix $\lambda\in\C^*$ that is not a root of unity.  Gluing the trivial line
bundle on $\P^1$ at $p$ and $q$ by the scalar $\lambda$ defines a line bundle
\(A_\lambda\in\Pic(C_0)\) with
\[
 \nu^*A_\lambda\simeq\cO_{\P^1},\qquad
 H^0(C_0,A_\lambda)=0,\qquad
 A_\lambda^{\otimes m}\not\simeq\cO_{C_0}\quad(m\ge1).
 \tag{4.4}
\]
\end{lemma}

\begin{proof}
For a nodal curve, a line bundle with trivial pullback to the normalization
is obtained by identifying the two fibres over the branches of the node.
Use the identification whose scalar is $\lambda$.  A global section pulls
back to a constant $c$ on the projective line and must satisfy
$c=\lambda c$; since $\lambda\ne1$, $c=0$.  The $m$th tensor power has
gluing scalar $\lambda^m$, which is never $1$ because $\lambda$ is not a
root of unity.  This proves (4.4).
\end{proof}

\begin{lemma}[an extending choice of gluing bundle]
\label{lem:extending-gluing}
After shrinking $\Delta$, the sections
\[
 O(t)=[0:1:0],\qquad
 P(t)=[3:\sqrt{36+t}:1]
\tag{4.5a}
\]
of $\mathscr C\to\Delta$ are disjoint and holomorphic.  If
$B=\cO_{\mathscr C}(P-O)$, then the line bundle
$\widetilde A=\operatorname{pr}_1^*B$ on $\mathcal X=\mathscr C\times\P^1$
restricts to $\operatorname{pr}_1^*A_{\lambda}$ on $S$ for
$\lambda=1/3$ (or its inverse, according to the gluing convention).  In
particular, the non-torsion bundle used in Theorem~\ref{thm:nodal-test} can
be chosen to extend to the displayed smooth total space.
\end{lemma}

\begin{proof}
The equation of $\mathscr C$ is satisfied by both sections; the square root
is holomorphic for $|t|$ small, and the two sections are disjoint because
$O$ has $z=0$ whereas $P$ has $z=1$.  On the normalization of $C_0$, write
$r=u/v$.  Then $O(0)$ corresponds to $r=\infty$, $P(0)$ to $r=2$, and the
two branches over the node correspond to $r=1$ and $r=-1$.  Hence
\[
 \nu^*(B|_{C_0})\simeq\cO_{\P^1}([2]-[\infty]).
\]
The rational function $f(r)=r-2$ has precisely this divisor and trivializes
the pullback.  Its values at the two branches are $f(1)=-1$ and
$f(-1)=-3$, so the induced gluing multiplier is $f(1)/f(-1)=1/3$,
up to inversion.  This is not a root of unity, and the lemma follows from
the gluing description in Lemma~\ref{lem:gluing-bundle}.
\end{proof}

\begin{theorem}[failure of the printed endpoint implication]
\label{thm:nodal-test}
Let $A=\operatorname{pr}_1^*A_\lambda$ on
$S=C_0\times\P^1$, with $\lambda$ as in Lemma~\ref{lem:gluing-bundle}.
Then
\[
 H^0(S,N^*_{S/\mathcal X}\otimes A)=0,
 \qquad
 H^0(S,\widetilde\Omega^1_S\otimes A)=0,
 \tag{4.5}
\]
but
\[
 H^0(S,\Omega^1_{\mathcal X}|_S\otimes A)\ne0.
 \tag{4.6}
\]
In particular, the two endpoint vanishings in (4.5) do not imply the
vanishing of (4.6), even for a proper morphism over a disc with smooth total
space and an irreducible reduced normal-crossing (self-crossing) central
fibre.
\end{theorem}

\begin{proof}
Because $S$ is the principal fibre $(t=0)$, its conormal line is trivial:
\(N^*_{S/\mathcal X}\simeq\cO_S\).  The projection
\(\operatorname{pr}_1:S\to C_0\) has
\(\operatorname{pr}_{1*}\cO_S\simeq\cO_{C_0}\), so (4.4) gives
\[
 H^0(S,A)\simeq H^0(C_0,A_\lambda)=0.
 \tag{4.7}
\]
This proves the first vanishing in (4.5).

We next check the second.  At a normal-crossing (self-crossing) point with local equation
\(uv=0\), the pullback map
\(\Omega^1_S\to\eta_*\Omega^1_{\widetilde S}\) has kernel exactly
\(\mathscr T_S\).  Indeed, in
\(R=\C\{u,v,w\}/(uv)\) the relation is
\(v\,du+u\,dv=0\).  If
\(a\,du+b\,dv+c\,dw\) restricts to zero on both branches, then
\(a\in(v)\), \(b\in(u)\), and
\(c\in(u)\cap(v)=(uv)=0\).  Modulo the relation the class is therefore a
multiple of $\tau=v\,du=-u\,dv$.  This class is nonzero: if
\(v\,du=h(v\,du+u\,dv)\), then coefficient comparison gives
  \((1-h)v=0\) and $hu=0$.  Since
\(\operatorname{Ann}_R(v)=(u)\) and \(\operatorname{Ann}_R(u)=(v)\),
\(1-h\in(u)\) and \(h\in(v)\), which would imply
\(1=(1-h)+h\in(u)+(v)\), impossible because $(u)+(v)$ is proper.
Moreover $(u+v)\tau=0$.  Since $u+v$ is a non-zero-divisor in
$R$ (the associated primes are $(u)$ and $(v)$), $\tau$ is genuine
torsion.
Conversely, each summand of the normalization module is a module over one of
the domains $R/(u)$ and $R/(v)$.  A non-zero-divisor of the reduced ring $R$
avoids both minimal primes and therefore acts injectively on both summands;
the target is torsion-free.  Thus the kernel of the normalization map is
exactly the torsion subsheaf, and there is an injection
\[
 \widetilde\Omega^1_S\hookrightarrow\eta_*\Omega^1_{\P^1\times\P^1}.
 \tag{4.8}
\]
Since $\eta^*A\simeq\cO_{\P^1\times\P^1}$, the projection formula and (4.8)
give
\[
 H^0(S,\widetilde\Omega^1_S\otimes A)
 \hookrightarrow
 H^0(\P^1\times\P^1,\Omega^1_{\P^1\times\P^1})=0.
 \tag{4.9}
\]
The last equality follows from
\(\Omega^1_{\P^1\times\P^1}\simeq
\operatorname{pr}_1^*\Omega^1_{\P^1}\oplus
\operatorname{pr}_2^*\Omega^1_{\P^1}\) and
\(H^0(\P^1,\Omega^1_{\P^1})=0\).

Finally, along the conductor $D$ the pullback of $dg|_S$ is zero.  Locally
$g=h\,uv$ with $h$ a unit, so after restriction to $S$ one has
$i^*(dg)=h(v\,du+u\,dv)$; both coefficients vanish on the preimage of $D$.
Proposition~\ref{prop:conductor-differential}
therefore gives a nonzero section in (4.6).  This proves the theorem.
\end{proof}

\begin{remark}[what the theorem does and does not refute]
\label{rem:scope}
The bundle $A$ is non-torsion on $S$: restriction to
\(C_0\times\{r\}\) would otherwise make $A_\lambda$ torsion.  Thus the
example meets the non-torsion condition attached to the p.~692 reduction.
In fact Lemma~\ref{lem:extending-gluing} gives an extending choice of $A$ on
the displayed smooth total space.  If $\widetilde L:=\omega_{\mathcal X}\otimes
\widetilde A^{-1}$, then
$(\widetilde L^\vee\otimes\omega_{\mathcal X})|_S\simeq A$, so the example
also has the Serre-dual twist appearing in the CDP argument.  This family
nevertheless does not have
the global topological and algebraic-dimension hypotheses of CDP20
Theorem~2.2.  The conclusion is exactly the formal one stated in
Theorem~\ref{thm:nodal-test}: the two displayed endpoint vanishings do not
control the torsion-supported lifting obstruction $\delta_A$.
\end{remark}

\section{The corrected reading of the CDP step}

The published paper states at p.~692 that, for an irreducible reduced fibre,
the conormal sequence makes it sufficient (when the restriction of the line
bundle is non-torsion) to prove the quotient vanishing (1.1)
\cite{CDP2020}.  Propositions~\ref{prop:filtration} and
\ref{prop:criterion} give the exact replacement:
\[
\begin{gathered}
 H^0(S,\widetilde\Omega^1_S\otimes A)=0,
 \quad H^0(S,N^*_{S/X}\otimes A)=0,\\[-2pt]
 \Longrightarrow\quad
 H^0(S,\Omega^1_X|_S\otimes A)=\ker\delta_A.
\end{gathered}
\tag{5.1}
\]
The missing assertion is therefore not another endpoint vanishing but the
injectivity of the explicitly defined map $\delta_A$.  For a reducible
reduced fibre the same formula must be applied to the whole Cartier fibre,
not separately to its components.  For a nonreduced fibre there are further
nilpotent layers, and (5.1) is not a complete analysis of that case.

The test family above is a counterexample to the unrestricted three-sheaf
implication for an arbitrary fibre line bundle $A$, not to CDP20
Theorem~2.2 or to its singular-fibre proposition under its full hypotheses.
Even though $A$ extends over this local total space, the family is not a
compact CDP threefold; whether the global hypotheses force a different
argument to make $\delta_A$ injective remains a separate question.

\subsection{A derived global criterion that survives the correction}

For completeness, we record a derived reformulation of the part of the CDP
strategy that does not use the defective endpoint reduction
\cite[Prop.~2.4]{CDP2020}.  The following generic-vanishing lemma is
elementary but useful in isolating the exact extra hypothesis; no priority is
claimed for this reformulation.

\begin{lemma}[generic endpoint vanishing]
\label{lem:generic-endpoints}
Let $Y$ be a connected compact complex manifold of dimension $n$, let $D$ be
a nonzero effective Cartier divisor with $\cO_Y(D)\in\Pic^0(Y)$, and let
$B$ be a holomorphic vector bundle.  Let $\mathcal P\subset\Pic^0(Y)$ be a
connected irreducible reduced complex analytic parameter space carrying a
Poincar\'e line bundle on $Y\times\mathcal P$.  Assume that dualization
preserves $\mathcal P$, and suppose it
  contains points $P_0=\cO_Y(-mD)$ and $Q_0=\cO_Y(-\ell D)$ for some
  $m,\ell\ge1$ with
  $H^0(Y,B\otimes P_0)=0$ and
  $H^0(Y,B^\vee\otimes\omega_Y\otimes Q_0)=0$ (the stabilization argument
  supplies such exponents whenever the corresponding twists lie in
  $\mathcal P$).  Then there is a dense open subset
$U\subset\mathcal P$ such that, for every $P'\in U$,
\[
 H^0(Y,B\otimes P')=H^n(Y,B\otimes P')=0.
 \tag{5.2}
\]
\end{lemma}

\begin{proof}
Put
  \(V_k=H^0(Y,B\otimes\cO_Y(-kD))\), viewed inside $H^0(Y,B)$, and note
  that $V_{k+1}\subseteq V_k$.  The descending chain stabilizes.  If $s$ is
  in its stable value, then at a point of $D$ choose a local defining
  nonzerodivisor $g_D$ for the Cartier divisor and a local frame of $B$.
  Every coefficient of the germ of $s$ is divisible by $g_D^k$ for every
  $k$.  The Krull-intersection theorem in the regular local ring
  $\cO_{Y,y}$ gives $\bigcap_{k\ge1}g_D^k\cO_{Y,y}=0$, so the germ is zero at
  every point of $D$.  Hence $s$ vanishes on a nonempty open set and, by
  connectedness, on $Y$.  Thus $V_k=0$ for all sufficiently large $k$.

  The chosen point $P_0$ satisfies
  $H^0(Y,B\otimes P_0)=0$.  The Poincar\'e bundle gives, locally on
$\mathcal P$, a family to which Grauert's theorem applies; the local jump
  loci $\{P':h^0(Y,B\otimes P')\ge1\}$ glue to a closed analytic subset of
  $\mathcal P$ (Grauert's theorem; see \cite{GrRe,BaSt}).  It is proper because it misses $P_0$, and irreducibility of
$\mathcal P$ makes its complement dense open.  Applying the same stabilization
argument to $B^\vee\otimes\omega_Y$ and using the chosen point $Q_0$ gives
$H^0(Y,B^\vee\otimes\omega_Y\otimes Q_0)=0$; dualization stability puts
$Q_0^\vee$ in $\mathcal P$.  Serre duality then gives
$H^n(Y,B\otimes Q_0^\vee)=0$.  The corresponding top-degree jump locus is
therefore also proper with dense open complement.  Intersecting the two
dense open complements proves the claim.
\end{proof}

\begin{theorem}[derived global CDP criterion]
\label{thm:corrected-global}
Let $X$ be a connected compact complex threefold with
\[
 H^1(X;\Z)=H^2(X;\Z)=0,\qquad a(X)=1,
 \tag{5.3}
\]
and suppose its algebraic reduction is a holomorphic map $f:X\to C$ to a
smooth compact curve.  If there is a line bundle $L_0\in\Pic(X)$ with
\[
 H^2(X,T_X\otimes L_0)=0,
 \tag{5.4}
\]
then a dense open subset $U\subset\Pic(X)$ satisfies
\[
  H^i(X,T_X\otimes M)=0\quad(i\in\{0,2,3\},\ M\in U),
 \qquad
 \left\langle c_3(T_X),[X]\right\rangle\le0.
 \tag{5.5}
\]
\end{theorem}

\begin{proof}
The exponential sequence and (5.3) identify
\(\Pic(X)\) with the complex vector group $H^1(X,\cO_X)$; in particular it
is connected and irreducible and is stable under dualization.  Choose a Stein
cover $\{U_i\}$ of $X$ and a \v{C}ech basis
$[a^{(1)}_{ij}],\ldots,[a^{(q)}_{ij}]$ of $H^1(X,\cO_X)$ represented by additive
cocycles.  If $\lambda=(\lambda_1,\ldots,\lambda_q)$ are the corresponding
coordinates, the functions
$\exp(\sum_r\lambda_r a^{(r)}_{ij})$ satisfy the multiplicative cocycle
condition and hence define the holomorphic Poincar\'e family on
$X\times\Pic(X)\simeq X\times H^1(X,\cO_X)$.  Every
line bundle in it has vanishing integral
first Chern class.  The algebraic-reduction map is surjective; choose
$c\in C$ and let $D=f^*(c)$ denote the scheme-theoretic fibre divisor.  It is
a nonzero effective Cartier divisor and
\(\cO_X(D)=f^*\cO_C(c)\in\Pic^0(X)\), because $H^2(X;\Z)=0$.
  The stabilization argument in Lemma~\ref{lem:generic-endpoints}, applied
  to $T_X$ and $T_X^\vee\otimes\omega_X$, supplies sufficiently large
  negative fibre twists; all of them lie in the vector group $\Pic(X)$.
  Apply that lemma with $Y=X$, $B=T_X$, and $\mathcal P=\Pic(X)$.  The
  exponential identification supplies the required Poincar\'e family and all
  negative fibre twists, so the lemma gives simultaneous degree-zero and
  degree-three vanishing on a dense open set.

The jump locus
\[
 J_2=\{M\in\Pic(X):h^2(X,T_X\otimes M)\ge1\}
\]
is a closed analytic subset by upper semicontinuity.  It is proper by (5.4),
so its complement is dense open.  Intersecting the three dense open sets
proves the cohomology assertions in (5.5).  On this intersection,
\[
 \chi(X,T_X\otimes M)=-h^1(X,T_X\otimes M)\le0.
 \tag{5.6}
\]

It remains to evaluate the Euler characteristic.  From (5.3), the rational
groups $H^1(X;\Q)$ and $H^2(X;\Q)$ vanish.  Universal coefficients and
Poincar\'e duality then give $H^4(X;\Q)=0$.  Hence
  $c_1(T_X)=c_1(K_X)=0$ and $c_2(T_X)=0$ in rational cohomology, and
  $c_1(M)=0$ as well.  Since $X$ is
  compact complex, holomorphic Hirzebruch--Riemann--Roch applies \cite{Hir}.  The
degree-six component of
\(\operatorname{ch}(T_X\otimes M)\operatorname{td}(X)\) is consequently
\(\tfrac12c_3(T_X)\), so
\[
 \chi(X,T_X\otimes M)=
 \frac12\left\langle c_3(T_X),[X]\right\rangle.
 \tag{5.7}
\]
Combining (5.6) and (5.7) proves the final inequality.
\end{proof}

\begin{corollary}[the intact relative hypothesis]
\label{cor:relative}
In the setting of Theorem~\ref{thm:corrected-global}, it is enough to assume
that one line bundle $L$ satisfies
\[
 R^2f_*(T_X\otimes L)=0.
 \tag{5.8}
\]
\end{corollary}

\begin{proof}
Choose a positive-degree line bundle $A$ on $C$ and set
$E=T_X\otimes L$.  Projection formula gives
\(R^2f_*(E\otimes f^*A^m)=0\).  For $m\gg0$, Serre vanishing on the curve
gives
\(H^1(C,R^1f_*E\otimes A^m)=0\), while
\(H^2(C,f_*E\otimes A^m)=0\) because $C$ has dimension one.  In the Leray
spectral sequence the three terms of total degree two are therefore
\[
\begin{aligned}
 E_2^{2,0}&=H^2(C,f_*E\otimes A^m)=0,\\
 E_2^{1,1}&=H^1(C,R^1f_*E\otimes A^m)=0,\\
 E_2^{0,2}&=H^0(C,R^2f_*E\otimes A^m)=0.
\end{aligned}
\]
Thus
\[
 H^2(X,E\otimes f^*A^m)=0.
\]
Apply Theorem~\ref{thm:corrected-global} to
$L_0=L\otimes f^*A^m$.
\end{proof}

\section{A local correction for ordinary relative one-forms}
\label{sec:relative-one-forms}

Alp{\"o}ge's source uses normal-crossing (self-crossing) charts for its reduced
singular fibre and cyclic charts for its multiple fibres
\cite[Sec.~10.2, source PDF pp.~87--88]{S6Source}.  The next
calculation is entirely local: it does not assume that these charts glue to a
compact threefold, and it concerns ordinary relative one-forms ($p=1$), not
their exterior powers.

\begin{proposition}[crossing and multiple-fibre modules]
\label{prop:relative-local}
The symbols $R_2$ and $R_3$ denote copies of the same regular local ring,
with different maps from the one-dimensional base:
let $R_2=\C\{x,y,z\}$ with base coordinate $t=xy$, and let
$R_3=\C\{x,y,z\}$ with base coordinate $t=xyz$.  Then
\[
\begin{aligned}
 M_2&:=\Omega^1_{R_2/\C\{t\}}
 =\frac{R_2\,dx\oplus R_2\,dy\oplus R_2\,dz}
        {R_2\,(y\,dx+x\,dy)},\\[2pt]
 M_3&:=\Omega^1_{R_3/\C\{t\}}
 =\frac{R_3\,dx\oplus R_3\,dy\oplus R_3\,dz}
        {R_3\,(yz\,dx+xz\,dy+xy\,dz)}.
\end{aligned}
\tag{6.1}
\]
Both modules are torsion-free, but neither is locally free at the origin.
They embed in free modules by
\[
\begin{aligned}
 M_2&\hookrightarrow R_2e\oplus R_2dz,
 &dx&\longmapsto xe,&dy&\longmapsto-ye,&dz&\longmapsto dz,\\
 M_3&\hookrightarrow R_3e_1\oplus R_3e_2,
 &dx&\longmapsto xe_1,&dy&\longmapsto ye_2,&dz&\longmapsto-z(e_1+e_2),
\end{aligned}
\tag{6.2}
\]
with images
\[
 R_2dz\oplus (x,y)e,
 \qquad
 R_3xe_1+R_3ye_2+R_3z(e_1+e_2).
\tag{6.3}
\]
If $R_m=\C\{s,y_1,y_2\}$ and the base coordinate is $t=s^m$ with
$m\ge2$, then
\[
\Omega^1_{R_m/\C\{t\}}
\simeq R_m\,dy_1\oplus R_m\,dy_2
       \oplus R_m/(s^{m-1})\,ds,
\tag{6.4}
\]
and the final summand is exactly its torsion submodule.
\end{proposition}

\begin{proof}
Differentiating the three equations gives the presentations in (6.1).  For
$M_2$, the first map in (6.2) sends $(a,b,c)$ to $(ax-by,c)$.  If this is
zero, then $c=0$ and $ax=by$.  Since $x$ and $y$ are relatively prime in the
UFD $R_2$, there is an $h$ with $b=xh$ and $a=yh$; hence its kernel is exactly
the displayed relation.  For $M_3$, the second map sends $(a,b,c)$ to
$(ax-cz,by-cz)$.  If it is zero, relative primeness first gives $x\mid c$,
then $y\mid(c/x)$; writing $c=xyh$ gives
$(a,b,c)=h(yz,xz,xy)$, again exactly the defining relation.  Thus both maps
are injective, proving torsion-freeness and (6.3).

At the maximal ideal $(x,y,z)$ each defining relation has all coefficients in
that maximal ideal.  Consequently the quotient of either module by the
maximal ideal has dimension $3$ over $\C$, while its generic rank is $2$;
a locally free module of rank $2$ would need only $2$ generators.  This proves
non-local-freeness at the origin.  Finally, for $t=s^m$ one has
$dt=m s^{m-1}ds=0$ relatively.  The integer $m$ is a unit in $\C$, so the
decomposition (6.4) follows.  Its last summand is annihilated by
$s^{m-1}$, while the other two summands are free; after inverting $s$ it
vanishes.  It is therefore exactly the torsion submodule.
\end{proof}

\begin{remark}[scope of the local correction]
\label{rem:relative-local-scope}
The modules in (6.1) are modules on the smooth total-space chart, and their
torsion-freeness is asserted before restricting to a fibre.  They must be
distinguished both from the intrinsic sheaf $\Omega^1_S$ of the singular
fibre and from the restricted relative module
$M_i\otimes_{R_i}R_i/(t)$ ($i=2,3$): tensoring with the fibre ring can create
conductor-supported torsion, as the $\tau$ calculation in the proof of
Theorem~\ref{thm:nodal-test} shows.  If a local quotient is finite
surjective \'{e}tale over a chart in (6.4), relative differentials commute with
this base change and faithful flatness then reflects torsion of the pulled-back
module; a ramified quotient
requires a separate invariant-module calculation.  Thus the bare $t=xyz$
chart has a torsion-free but non-locally-free total-space relative module for
$p=1$; this fact alone does not decide torsion after fibre restriction or
quotient.  The multiple-fibre charts do have the total-space torsion in
(6.4).  No claim about torsion in $\bigwedge^2\Omega^1$ or
$\bigwedge^3\Omega^1$, and no global direct-image or Hodge conclusion, is
made here.
\end{remark}

\section{\texorpdfstring{Relation to the proposed $S^6$ reconstruction}{Relation to the proposed S6 reconstruction}}

The larger validate/repair/disprove workbench applies the filtration above to
a proposed compact complex threefold with a non-normal singular fibre and
records candidate values
\[
 H^1(X;\Z)=H^2(X;\Z)=H^3(X;\Z)=0,
 \qquad a(X)=1,
 \qquad \langle c_3(T_X),[X]\rangle=2.
 \tag{7.1}
\]
It also records candidate-specific calculations of integral specialization,
Leray transgressions, and the nonvanishing of
$R^2f_*(T_X\otimes L)$.  Those calculations are the source manuscript's
application and are not repeated here.  None of the upstream construction,
topology, or Picard assertions is proved in this short note.  Applying Proposition
\ref{prop:criterion} to that candidate would require, at minimum, a complete
independent verification of the smooth compact threefold, the full reduced
Cartier fibre, the normalization and line-bundle restrictions, and the
actual connecting map $\delta_A$ (or an equivalent explicit computation).
The present result supplies the local algebraic target for that check; it is
not an adjudication of the candidate's existence.

\section*{AI-use and authorship disclosure}

Generative-AI systems assisted the transcription of source material, symbolic
and exact-arithmetic checks, literature searching, and preparation of this
draft.  The local drafting and audit assistance used OpenAI Codex gpt-5.6-sol,
Ultra; no AI system is an author.  A human author must verify every statement,
secure the right to submit the work, and take responsibility for the final
version, consistent with the current arXiv policy for generative-AI tools
\cite{arXivModeration}.

\bibliographystyle{amsplain}
\bibliography{references}

\end{document}
